% Insert after the paper defines pair gate G and source timestamps.
% Required packages: amsmath, amssymb, amsthm.
% Required environments: theorem, proposition, corollary, proof.

\paragraph{Direct temporal witnesses.}
Let $G(d,c)$ test whether claim $d$ supplies a replacement for claim $c$.
The test uses claim text and fixed extraction rules.
It need not define a transitive relation.
At cutoff $\tau$, eligible evidence is
\begin{equation}
 A_G(\tau)=\{c:t_c\leq\tau,\;\nexists d\text{ with }t_c<t_d\leq\tau\text{ and }G(d,c)\}.
 \label{eq:direct-eligibility}
\end{equation}
This definition concerns the supplied corpus and its timestamp semantics.
Equal timestamps do not establish replacement order.

For each claim, store its first direct replacement time:
\begin{equation}
 b_G(c)=\min\{t_d:t_d>t_c,\;G(d,c)\},\qquad
 I_G(c)=[t_c,b_G(c)),
 \label{eq:direct-certificate}
\end{equation}
where $\min\varnothing=\infty$.
Store one attaining witness for source attribution.
An identifier rule resolves ties without changing the boundary.

\begin{theorem}[All-cutoff compilation]
\label{thm:all-cutoff}
For every finite cutoff $\tau$, $c\in A_G(\tau)$ if and only if $\tau\in I_G(c)$.
The compiler therefore preserves all temporal masks with at most one witness per claim.
\end{theorem}
\begin{proof}
If $t_c>\tau$, both statements are false.
Otherwise, a qualifying witness exists by $\tau$ exactly when $b_G(c)\leq\tau$.
Taking the complement proves the claim.
\end{proof}

\begin{corollary}[Prefix consistency]
\label{cor:prefix}
Querying the full index at $\tau$ equals compiling only claims dated at or before $\tau$.
Consequently, later additions cannot change earlier evidence states.
\end{corollary}
\begin{proof}
Both procedures test exactly the pairs satisfying $t_c<t_d\leq\tau$.
\end{proof}

\paragraph{Local influence.}
The guarantee does not require error-free matching.
It controls how matching errors affect the compiled evidence.

\begin{theorem}[Bounded influence]
\label{thm:influence}
Fix timestamps and a candidate ranking.
Suppose gates $G$ and $H$ differ on $h$ temporally ordered pairs.
Let $U$ contain their older endpoints.
For every cutoff $\tau$,
\begin{align}
 A_G(\tau)\triangle A_H(\tau)&\subseteq U,\label{eq:mask-influence}\\
 |R_G(\tau)\triangle R_H(\tau)|&\leq2\min\{k,|U|\}\leq2h,
 \label{eq:context-influence}
\end{align}
where $R$ filters the fixed ranking before selecting at most $k$ claims.
\end{theorem}
\begin{proof}
The boundary for $c$ depends only on pairs whose older endpoint is $c$.
It cannot change outside $U$.
Changing one eligibility bit replaces at most one selected claim.
Applying this argument successively gives $2|U|$.
The two contexts contain at most $2k$ claims altogether.
\end{proof}

\paragraph{Why direct witnesses matter.}
Connected components treat local matches as a transitive identity relation.
One wrong bridge can therefore change many historical intervals.
Consider repeated value $a$ for one key at times $0,\ldots,m-1$.
It changes to $b$ at time $2m$.
A second key contains $z$ at $m$ and $w$ at $2m+1$.
Let all four values differ.
A false bridge between $b$ and $w$ merges the histories.
Component sorting then closes all $m$ earlier $a$ claims at time $m$.
At cutoff $m+\tfrac12$, direct compilation changes no claim.
The false bridge affects only $b$, which has not arrived.
Thus, component closure can amplify one future error into arbitrarily many earlier exclusions.

\paragraph{Construction cost.}
Process timestamps in increasing batches.
Keep unresolved older claims in a complete inverted index.
The first positive gate fixes a claim's boundary and removes it from the index.
Insert each batch after all comparisons, so tied claims remain unordered.
The output uses $O(n)$ words.
Construction costs $O(n\log n+L+V+Pg)$ after feature extraction.
Here $L$ counts posting entries, $V$ counts posting visits, and $P$ counts pair tests.
The quantity $g$ bounds one gate evaluation.
Uninformative blocks can still require quadratic work.
